\chapter{Contratos entre as etapas}
\label{ap:contratos}

% Fonte: project-src/contracts/README.md e os esquemas JSON.

Os quatro contratos da Tabela~\ref{tab:contratos} foram escritos como esquemas
JSON, com exemplos, e não como texto. Cada documento carrega uma versão; um
consumidor aceita qualquer documento com a mesma versão principal e ignora
campos desconhecidos, e mudanças incompatíveis exigem nova versão principal.

\section{C1: arquivo de configuração}

Descreve o modelo de interrupções de um programa: os fluxos, com ponto de
entrada e prioridade; as primitivas de mascaramento; o tratamento de funções
externas; e as opções semânticas, como a preempção entre fluxos de mesma
prioridade e a reentrância. As opções que afetam o recall trazem o valor
conservador como padrão.

\section{C2: registro de contexto}

Um documento por candidato, com os campos da Tabela~\ref{tab:registro}. Cinco
decisões de projeto estão codificadas na própria estrutura, para que não possam
ser revertidas sem que alguém perceba:

\begin{enumerate}
  \item o mascaramento tem três valores (desabilitado, habilitado e desconhecido), e não existe um campo booleano ``mascarado'';
  \item o estado de mascaramento num ponto e ao longo de um intervalo são campos
    distintos, pois responder um e inferir o outro perde defeitos;
  \item o veredicto do solver tem quatro valores, e o inconclusivo exige um
    motivo;
  \item a proveniência dos fatos é obrigatória;
  \item não existe campo de fatiamento: o contexto adicional é obtido pelo
    contrato C4.
\end{enumerate}

Cada registro traz ainda uma impressão digital que depende apenas do programa,
da classe, da variável e das posições dos acessos, usada como chave do cache de
respostas do modelo.

\section{C3: organização de uma execução}

Define os diretórios e arquivos de uma execução da ferramenta: um manifesto com
a configuração, as versões e as decisões de avaliação usadas; os candidatos de
cada estágio; e o histórico das conversas com o modelo. Uma execução em que algum
candidato não tenha nem registro de contexto nem prova de impossibilidade é
rejeitada pela verificação do próprio formato.

\section{C4: pedidos de contexto}

Define o que o modelo pode pedir e como a resposta volta. Os tipos de pedido
incluem a definição de uma função, o corpo de uma \isr, o estado de mascaramento
numa linha, a expansão de uma macro, todos os acessos a uma variável e todos os
caminhos de chamada até uma função. Toda resposta traz um estado (resolvido, não encontrado, não suportado, truncado, ambíguo ou recusado) e, quando o
pedido não pôde ser atendido, a regra de recall que se aplica. Uma definição
não encontrada, por exemplo, deve ser tratada como uma função que pode acessar
a variável; uma resposta vazia nunca pode ser lida como ``não há nada''.
